% =====================================================================
% Section 6 — Conclusion
% =====================================================================

\section{Conclusion}\label{sec:conclusion}

This work proposed a quantitative-coupled BowTie risk analysis of the
valve train of a high-performance cylinder head conceived for the
Volkswagen AP 1.8 engine operating at 10\,000~rpm. The qualitative
BowTie diagram organises six threats and five preventive barriers
around the top event \emph{loss of follower contact between valve and
cam}, together with five consequences and five mitigative barriers. To
overcome the qualitative limitation of the classical BowTie, each
preventive barrier was associated with a deterministic indicator
(dynamic follower-contact ratio, spring natural frequency,
coil-bind margin, Goodman fatigue factor of safety, and a
closed-form differential hot-clearance estimate), all anchored in
literature-based pass criteria. A probabilistic refinement
propagated the catalogued ASTM~A877/A877M wire tolerances and
strength dispersion through the indicators using a Monte Carlo
simulation with $N=2\times 10^{5}$ trials, converting nominal-point
status labels into manufacturing-realistic probabilities of barrier
compliance.

For the AP 1.8 head under the operating envelope adopted, the analysis
identified two concurrent at-risk preventive barriers in the baseline
configuration and one additional marginal barrier exposed by the
probabilistic refinement:
(i) the \emph{outer-spring coil-bind margin}
($M_{\mathrm{out}}=0.78$~mm at nominal, against the project criterion
$M\geq 1.5$~mm, with $P[\mathrm{fail}]=77.9\,\%$ under ASTM~A877
tolerance);
(ii) the \emph{inner-spring fatigue Goodman factor of safety}
($n_{f,\mathrm{in}}\approx 0.68$ at nominal, against $n_{f}\geq 1$,
with $P[n_{f}<1]=100\,\%$ across the manufactured population); and
(iii) the \emph{inner-spring coil-bind margin}, classified as
active at the nominal point ($M_{\mathrm{in}}=2.38$~mm) but
re-classified as marginal under tolerance propagation
($P[\mathrm{fail}]=16.2\,\%$), a finding that the deterministic
nominal-point analysis does not surface. The two at-risk barriers
affect different physical components and require independent design
interventions. The \emph{outer-spring natural frequency}
($f_{n,\mathrm{out}}/f_{\mathrm{cam,max}}\approx 3.81$,
$P[\mathrm{fail}]=94.4\,\%$ against the conventional 4$\times$
rule) emerges as an additional residual concern. The remaining
barriers --- dynamic follower contact
($R\geq 1.343$ in the worst case at 10\,000~rpm; the
6\,000--12\,000~rpm sweep of Section~\ref{sec:res-envelope}
identifies the loss-of-contact boundary at approximately
11\,500~rpm), outer-spring fatigue ($n_{f,\mathrm{out}}\approx 1.27$,
$P[\mathrm{fail}]<0.01\,\%$), inner-spring surge ratio, and the
differential hot-clearance indicator
($c_{\mathrm{cold,req}}=1.25$~mm, bracket [1.13, 1.60]~mm) ---
satisfy the pass criteria with non-uniform margins. The comparison
with a conventional FMEA scored under the AIAG--VDA 2019 rubric
showed that the quantitative-coupled BowTie offers (i) explicit
visualisation of the causal chain, (ii) reproducibility from the
sizing calculations, and (iii) identification of multiple concurrent
critical barriers via a sensitivity analysis that the harmonised
ordinal score would have missed.

The proposed methodology is reproducible and can be transposed to
other high-performance cylinder heads operating at high rotational
speed, provided the analytical sizing of the valve train is
available. A construct-validity benchmark against four documented
valve-train and engine-component failure case studies from the recent
literature \citep{velezgodino2019overheadvalvetrain,soffritti2018wornvalvetrain,xing2021valvespringEAC,zhao2023camshaftV6}
shows that three of the four cases are accounted for within the
present framework as formulated --- one is caught quantitatively by
the deterministic indicator set (Xing, via the Goodman factor of
safety), one is caught on the mitigative side as designed
(V\'elez Godi\~no), and one is flagged qualitatively and motivates a
candidate assembly-tolerance preventive barrier (Soffritti) --- while
the fourth case (Zhao) defines the boundary of the current barrier
set and motivates a documented lubrication-regime extension. The
benchmark confirms construct validity but does not establish
predictive validity nor the absolute magnitude of the indicator
margins for the AP 1.8 head; a test-rig campaign with statistically
specified acceptance criteria therefore remains necessary, and its
scope is specified next.

\paragraph{Validation roadmap with statistical acceptance criteria.}
The minimum experimental verification programme that would close the
empirical gap on the analytical findings for the AP 1.8 head
comprises three coordinated stages. Each stage is specified with an
explicit sample size, instrumentation envelope, and statistical
acceptance rule keyed to the analytical predictions of
Section~\ref{sec:res-quant} and the probabilistic envelope of
Section~\ref{sec:res-mc}.

\begin{enumerate}
\item \textbf{Spring rig --- modal verification ($n\geq 3$
specimens per spring).} A single-spring test fixture
instrumented with a non-contacting laser-displacement sensor
($\geq$ 1~kHz sampling) on the active coils identifies the surge
response across the 50--600~Hz band, under the catalogued
seat-preload boundary condition.
\emph{Acceptance criterion}: the sample-mean first surge frequency
$\bar{f}_{n}$ from $n\geq 3$ specimens shall fall within the 90\,\%
confidence interval of the analytical prediction (calibrated to the
MC distribution of Section~\ref{sec:res-mc}, i.e.\
$f_{n,\mathrm{out}}\in [304,\,334]$~Hz and
$f_{n,\mathrm{in}}\in [545,\,600]$~Hz at 90\,\% confidence). A
failure of containment is reported as a falsification of the
analytical model and triggers re-evaluation of the active-mass
hypothesis (total versus Rayleigh-effective).

\item \textbf{Motored-engine speed sweep --- coil-bind and
follower-contact verification.} A motored-engine campaign from
6\,000 to 12\,000~rpm engine speed (the dynamic-loss-of-contact
boundary identified in Section~\ref{sec:res-envelope}) is
conducted with a non-contacting valve-position transducer sampled
at $\geq$~50~kHz to resolve the deceleration peak of the cam
profile. Two acceptance criteria are applied:
\emph{(2a)} absence of valve float at the deceleration peak
($a_{\max,-}=20\,362$~m\,s$^{-2}$, Section~\ref{sec:res-quant})
verified by zero-crossing analysis of the
(measured$-$predicted) lift residual at each rpm point, with a
detection threshold of $\geq$~0.05~mm departure for $\geq$~3
consecutive cam events;
\emph{(2b)} the residual coil-bind margin measured at full lift
($n\geq 5$ measurements at 10\,000~rpm) shall lie within
$\pm$0.1~mm of the MC 50\,\% credible interval
$M_{\mathrm{out}}\in [0.51,\,1.10]$~mm for the baseline outer
spring (or $M_{\mathrm{out}}\approx 2.0$~mm post-correction). A
single instance of $M_{\mathrm{out}}<0$~mm (hard contact) at any
rpm within the safe envelope falsifies the corrected design.

\item \textbf{Inner-spring fatigue campaign --- Goodman criterion
verification ($n\geq 5$ specimens).} A representative sample of
inner springs is cycled under shear-stress amplitudes reproducing
the analytical
$\tau_{m,\mathrm{in}}=858$~MPa and
$\tau_{a,\mathrm{in}}=374$~MPa, with each specimen run to either
run-out at $10^{8}$ cycles or first-failure indication.
\emph{Acceptance criterion}: with $n=5$ specimens and assuming a
Weibull failure model, the upper one-sided 90\,\% confidence
bound on the empirical Goodman factor of safety
$\hat{n}_{f,\mathrm{in}}$ shall remain below unity ($\leq 0.78$
upper bound for 0 of 5 specimens reaching run-out), thereby
confirming the analytical prediction
$n_{f,\mathrm{in}}\approx 0.68$ and the MC 5\,\% quantile
$n_{f,\mathrm{in,5\%}}=0.637$. A pass condition
($\hat{n}_{f}\geq 1$ upper bound) on 5 of 5 specimens reaching
run-out would falsify the predicted Goodman deficit and require
re-evaluation of the shear endurance limit calibration.
\end{enumerate}

The total test-rig hours for this minimum programme are estimated at
the order of 600--800~h. Subsequent extensions of the work will
address (i) a second-generation Monte Carlo propagation that adds
correlated operational sources (thermal-soak transients, oil
condition, valve-seat wear) and joint barrier-failure probabilities;
(ii) extension of the preventive set with candidate barriers on
assembly tolerance and on the lubrication regime motivated by the
construct-validity benchmark of Section~\ref{sec:res-litvalidation};
(iii) coupled thermal--structural finite-element analysis of the
cylinder head to refine the differential hot-clearance estimate of
Section~\ref{sec:res-quant} and tighten the
$c_{\mathrm{cold,req}}\in [1.13,\,1.60]$~mm bracket; and
(iv) integration of the resulting indicator set as a constraint in
the multi-objective optimisation of the cylinder-head geometry.
